
\begingroup
\renewcommand{\thefootnote}{*}%
\section*{6. Subfunctors and group subschemes%
  \texorpdfstring{\protect\footnotemark}{}}
\footnotetext{\label{note:sga3-VIB-section6-star}On the same
theme, see also the results of Expos\'e~XI, \S6, whose natural
place would be in the present Expos\'e~VI\(_B\). There one finds,
in particular, representability criteria for certain group
subfunctors of a given group scheme.}
\endgroup
\phantomsection\label{sec:sga3-VIB-6}%
\addcontentsline{toc}{section}{6. Subfunctors and group subschemes}

\medskip
\phantomsection\label{def:sga3-VIB-6.1}%
\noindent\textbf{Definition 6.1.---}
\textup{(i)} Let \(X\) be an \(S\)-functor (that is, a functor
from \((\mathbf{Sch}/S)^\circ\) to \((\mathbf{Ens})\)), let \(G\)
be an \(S\)-functor in groups, and let \(u\) and \(v\) be two
\(S\)-morphisms from \(X\) to \(G\). The \emph{transporter of
\(u\) into \(v\)}, denoted
\(\underline{\operatorname{Transp}}(u,v)\), is the
sub-\(S\)-functor \(T\) of \(G\) defined by
\[
  T(S')=\bigl\{g\in G(S')\bigm|(\operatorname{int}g)\circ u_{S'}=v_{S'}\bigr\}.
\]

\smallskip
\noindent\textup{(ii)} Let \(G\) be an \(S\)-functor in groups,
and let \(X\) and \(Y\) be two sub-\(S\)-functors of \(G\). The
\emph{transporter of \(X\) into \(Y\)} (respectively, the
\emph{strict transporter of \(X\) into \(Y\)}), denoted
\(\underline{\operatorname{Transp}}_G(X,Y)\) (respectively,
\(\underline{\operatorname{Transp}}^{\mathrm{str}}_G(X,Y)\)),
is the sub-\(S\)-functor \(T\) of \(G\) defined by
\[
  T(S')=\bigl\{g\in G(S')\bigm|(\operatorname{int}g)(X_{S'})\subset Y_{S'}\
  \bigl(\text{respectively, }(\operatorname{int}g)(X_{S'})=Y_{S'}\bigr)\bigr\}.
\]
Note that
\phantomsection\label{eq:sga3-VIB-6.1-inversion}%
\[
\underline{\operatorname{Transp}}^{\mathrm{str}}_G(X,Y)
 =
\underline{\operatorname{Transp}}_G(X,Y)
 \cap
c\bigl(\underline{\operatorname{Transp}}_G(Y,X)\bigr),
\]
where \(c\) denotes the symmetry of \(G\).%
\footnote{Translator's note: the Gille--Polo re-edition (2011), Expos\'e~VI\textsubscript{B}, observes
that (ii) also serves for arbitrary morphisms of \(S\)-functors \(u\colon X\to G\) and \(v\colon Y\to G\), through
their images, the subfunctors \(u(X)\colon S'\mapsto u(S')(X(S'))\) and \(v(Y)\) of \(G\): one may form
\(\underline{\operatorname{Transp}}_G(u(X),v(Y))\), which is \(\underline{\operatorname{Transp}}_G(X,Y)\) when
\(u\) and \(v\) are the inclusions of subfunctors; it states 6.2 below, and 10.11, in these terms.}

\smallskip
\noindent\textup{(iii)} Let \(G\) be an \(S\)-functor in groups,
let \(X\) be an \(S\)-functor, and let \(u\) be an \(S\)-morphism from
\(X\) to \(G\). The \emph{centralizer of \(u\)} is the group
sub-\(S\)-functor
\(\underline{\operatorname{Centr}}(u)=\underline{\operatorname{Transp}}(u,u)\)
of \(G\).

\smallskip
\noindent\textup{(iv)} Let \(G\) be an \(S\)-functor in groups,
let \(H\) be a sub-\(S\)-functor of \(G\), and let
\(i:H\to G\) be the canonical \(S\)-morphism. The
\emph{centralizer} (respectively, the \emph{normalizer}) \emph{of \(H\) in
\(G\)} is the group sub-\(S\)-functor of \(G\)
\phantomsection\label{eq:sga3-VIB-6.1-centralizer-normalizer}%
\[
\underline{\operatorname{Centr}}_G H
  =\underline{\operatorname{Centr}}(i)
  =\underline{\operatorname{Transp}}(i,i)
\qquad
\bigl(\text{respectively, }\underline{\operatorname{Norm}}_G H
  =\underline{\operatorname{Transp}}^{\mathrm{str}}_G(H,H)\bigr).
\]
Finally, the \emph{center of \(G\)} is the \(S\)-functor in
groups
\(\underline{\operatorname{Centr}}\,G=\underline{\operatorname{Centr}}_G G\).\footnote{Translator's note: the
Gille--Polo re-edition (2011), Expos\'e~VI\textsubscript{B}, adds that these functors commute with base change: if
\(G',X',Y',u',v'\) are deduced from \(G,X,Y,u,v\) by a base change \(S'\to S\), then
\(\underline{\operatorname{Transp}}(u,v)_{S'}=\underline{\operatorname{Transp}}(u',v')\) and
\(\underline{\operatorname{Transp}}_G(X,Y)_{S'}=\underline{\operatorname{Transp}}_{G'}(X',Y')\), and similarly for
the strict transporter, centralizers and normalizers; indeed, for an \(S'\)-scheme \(T\), both sides have the same
\(T\)-valued points, the defining conditions involving only \(T\) and the \(S\)-schemes over it.}

